\section{Background}

The security evidence for the Android app encompasses a wide range of features, including manifest declarations, requested permissions, exported components, packaged resources, embedded libraries, native code, storage configuration, and network-security settings. Android permissions protect restricted data and actions \cite{androidPermissions2026}. The Android privacy guide encourages minimizing data access \cite{androidPrivacyChecklist2026}, and the Android security guide addresses secure inter-app communication and related platform practices \cite{androidSecurityBestPractices2026}. These static surfaces are useful indicators, but do not prove runtime behavior.

Static analysis is repeatable and versionable: it inspects the packaged APK evidence before execution and can be linked to the exact package/version identity. Detector coverage and reporting consistency still vary between Android SAST tools, so the findings in this paper are treated as exposure indicators, not as proof of exploitability. Runtime analysis complements this evidence under controlled conditions. Packet captures can support traffic-shape, DNS, TLS-fingerprint, and protocol-prevalence measurements, but a non-root setting without MITM interception does not reveal encrypted payload semantics. We therefore report runtime metadata and derived indicators only.

The modern Android distribution also affects reproducibility. Google Play commonly delivers device-specific APKs sets from Android App Bundles, including base and split APKs \cite{androidAppBundles2026}. A reproducible study must retain package name, version code/name, APK hashes, split membership, static run identifiers, dynamic run identifiers, capture timestamps, and quality-assurance state. This build-backed linkage is the foundation for the static runtime alignment used below.
